\section{Mapa de Navegación}

Como destaca Spool (2005), cuando un usuario entiende la información y cómo está organizada, se siente confiado y aumenta su nivel de \textit{entusiasmo} al usar el sistema. Construimos el mapa de navegación siguiendo los 6 pasos que se describen a continuación:

\subsection{Revisión de los Resultados de la Arquitectura}
Iniciamos la fase de diseño con una revisión de la Arquitectura de Información propuesta en la sección anterior, ya corregida con la evidencia de la prueba de árbol. El antiguo menú lateral de ``Contenido temático'' del portal CIF repartía el acervo en más de 20 carpetas alfabéticas y generaba una fuerte sobrecarga visual y cognitiva. Frente a ese punto de partida, determinamos que la navegación principal de Karisma Data se consolidaría en solo cuatro módulos base: \textbf{Inicio}, \textbf{Exploración y extracción}, \textbf{Gobierno del dato} y \textbf{Administración}.

El primer nivel resistió la prueba y se conservó sin cambios: cuatro módulos bastaron para cubrir las 35 tarjetas del instrumento y los ocho perfiles de la Actividad 1, frente a las más de 20 entradas alfabéticas del portal legado. Cambiaron el nombre de la tercera categoría, que dejó de llamarse ``Gobierno y Supervisión'' para llamarse \textbf{Gobierno del dato}, y el contenido de la segunda, que absorbió los tableros e indicadores de riesgo. El paso 4 documenta y justifica ambos ajustes.

\subsection{Determinación de Subcategorías}
Dentro de cada categoría principal identificamos y organizamos las subcategorías para profundizar en el contenido sin sobrecargar al lector. El segundo nivel quedó así:

\begin{uxlist}
  \item \textbf{1. Inicio:} buscador unificado, búsquedas recientes, favoritos, mis alertas y mi perfil.
  \item \textbf{2. Exploración y extracción:} 2.1 Catálogo temático (cartera de crédito, liquidez, derivados, captación, mercado de dinero, acciones); 2.2 Consulta y filtros (buscador facetado, filtros avanzados, vista previa, compartir consulta); 2.3 Exportaciones (exportación CSV/Excel, historial de exportaciones); \textbf{2.4 Tableros e indicadores} (tablero ejecutivo, indicadores de riesgo, señales de riesgo).
  \item \textbf{3. Gobierno del dato:} 3.1 Diccionario y metadatos (diccionario de datos, metadatos, reglas de negocio); 3.2 Linaje y calidad (linaje, calidad de datos, monitoreo de cargas); 3.3 Catálogo de fuentes, con faceta transversal.
  \item \textbf{4. Administración:} 4.1 Usuarios, roles y permisos; 4.2 Solicitudes y aprobaciones (solicitar acceso, estado de solicitudes, flujos de aprobación); 4.3 Bitácora de accesos, con acceso cruzado desde Gobierno del dato; 4.4 Integraciones (credenciales API, consumo por API, conectores de fuentes).
\end{uxlist}

La subcategoría 2.4 es la única rama incorporada después de la prueba y su origen se explica en el paso 4. Dentro de ``Exploración y extracción'' definimos además un modelo de navegación transversal que abarca todos los temas operativos de la institución mediante un esquema de búsqueda por facetas. Ese esquema permite el etiquetado múltiple de las fuentes que impactan a distintas áreas de forma simultánea, evita la duplicación de archivos y elimina la duda de ``en qué carpeta exacta guardaron este reporte''. Nueve tarjetas quedaron marcadas como facetas transversales: consumo por API, calidad de datos, vista previa, catálogo de fuentes, mis alertas, credenciales API, monitoreo de cargas, permisos e historial de exportaciones.

\subsection{Visualización de la Jerarquía de Navegación}
Con las jerarquías claras, generamos una estructura visual (Mapa de Sitio) que representa todas las rutas posibles dentro de la aplicación. La Figura \ref{fig:mapa} corresponde a la versión revisada e incorpora los cuatro cambios documentados en el paso 4: los tableros e indicadores cuelgan de Exploración y extracción como rama 2.4, la tercera categoría aparece rotulada ``Gobierno del dato'', la bitácora de accesos vive en Administración con un acceso cruzado desde Gobierno del dato, y las facetas transversales quedan marcadas sobre las tarjetas que las recibieron.

\figuraux{figuras/a3_mapa_sitio.png}{Estructura jerárquica (Mapa de Navegación) del portal institucional, en su versión revisada tras la prueba de árbol.}{fig:mapa}

Como se observa en la Figura \ref{fig:mapa}, el mapa implementa el principio de \textbf{revelación progresiva}: la barra lateral mantiene visibles los cuatro módulos, pero despliega el segundo nivel únicamente del módulo en uso. El perfil operativo (Laura Méndez) encuentra su barra de búsqueda en ``Inicio'' de inmediato, sin tropezar visualmente con las herramientas de ``Administración'' que necesita el perfil de administración (Mariana Ovalle Ríos), y el perfil directivo (Arturo Castañeda) alcanza el tablero ejecutivo dentro del mismo módulo donde ya consulta las cifras, sin cruzar hacia gobierno. Ninguna ruta termina en un callejón sin salida: cada hoja del árbol ofrece al menos una salida hacia una tarea contigua, ya sea exportar el resultado, abrir su ficha de metadatos o solicitar el acceso que falta. Las facetas transversales refuerzan esa propiedad, porque una misma tarjeta resulta alcanzable desde más de una rama sin que el contenido se duplique.

\subsection{Validación y Mejora}
El mapa se sometió a prueba antes de darlo por bueno. El 9 de agosto de 2026, al cierre de cada una de las ocho sesiones de \textit{card sorting} abierto, aplicamos la prueba de árbol de cinco tareas a los ocho evaluadores prototipo condicionados por las personas de la Actividad 1, sin mostrarles interfaz alguna y con las sesiones a ciegas entre sí. Cada evaluador eligió un módulo de primer nivel por tarea, lo que produjo 40 observaciones; 28 de ellas coincidieron con la ruta esperada. A esa fuente se sumó una segunda, recogida durante el propio ejercicio de clasificación: las peticiones de colocar una misma tarjeta en dos o más grupos a la vez. Los 8 evaluadores pidieron al menos cinco duplicados cada uno.

El desempeño por tarea fue desigual. Exportar la cartera de crédito resultó unánime: los 8 evaluadores fueron a Exploración. Crear una credencial de integración lo resolvieron 7 de los 8 en Administración. Localizar una fuente de liquidez y revisar quién autorizó un acceso los resolvieron 6 de los 8 cada una. Consultar un indicador de riesgo, en cambio, lo acertó \textbf{solo 1 de los 8 evaluadores}: 6 lo buscaron en Exploración y 1 en Inicio, y titubearon 6 antes de decidir. Ese es el fallo principal de la arquitectura propuesta y obligó a mover contenido de módulo. La segunda señal fue más silenciosa: en la tarea de revisar quién autorizó un acceso titubearon los 8 evaluadores, incluidos los 6 que terminaron acertando, y 2 de ellos se fueron a Gobierno.

Los comentarios recibidos se incorporaron al mapa y quedaron registrados en la tabla de cambios, con la versión anterior, la observación recibida, la modificación aplicada y su justificación. Así el refinamiento se puede rastrear sin confundir decisiones internas del equipo con evidencia de la prueba.

\begin{uxtabla}{L{2.4cm}Z{1.05}Z{1.0}Z{0.95}}{Tabla de cambios del mapa de navegación tras la prueba de árbol y el ejercicio de clasificación}
  \uxheadrow \thd{Versión anterior} & \thd{Observación recibida} & \thd{Modificación aplicada} & \thd{Justificación} \\
  Tableros e indicadores de riesgo colgaban de ``Gobierno y Supervisión'' & Solo 1 de los 8 evaluadores los buscó ahí; 6 de los 8 fueron a Exploración y 1 a Inicio, con 6 titubeos & Se movieron a Exploración y extracción como subcategoría ``2.4 Tableros e indicadores'': tablero ejecutivo, indicadores de riesgo y señales de riesgo & Consultar un indicador es una tarea de consulta, no de gobierno; la matriz de similitud agrupa esas tres tarjetas en un conglomerado propio, separado del bloque de definición del dato \\
  La tercera categoría se llamaba ``Gobierno y Supervisión'' & Al salir los tableros, la categoría ya no supervisa nada: conserva únicamente el marco del dato & Se renombró ``Gobierno del dato'' y agrupa diccionario, metadatos, reglas de negocio, linaje, calidad y monitoreo de cargas & La etiqueta debe describir lo que el módulo contiene; un rótulo que promete supervisión y entrega definiciones reintroduce la desorientación que la arquitectura buscaba eliminar \\
  ``Bitácora de accesos'' existía solo en Administración & Titubearon los 8 evaluadores en la tarea de revisar quién autorizó un acceso y 2 de los 8 la buscaron en Gobierno & Se conservó en Administración como rama 4.3 y se añadió un acceso cruzado desde Gobierno del dato & La matriz la agrupa con fuerza junto a usuarios, permisos y solicitudes, pero los perfiles de auditoría y de propietario del dato la buscan desde gobierno: mover la tarjeta habría roto un conglomerado sólido, duplicar la ruta no \\
  Jerarquía de pertenencia única: cada tarjeta vivía en un solo nodo & Los 8 evaluadores pidieron al menos cinco duplicados cada uno; consumo por API y calidad de datos los pidieron 6 de los 8 & Se definieron nueve facetas transversales de etiquetado múltiple: consumo por API, calidad de datos, vista previa, catálogo de fuentes, mis alertas, credenciales API, monitoreo de cargas, permisos e historial de exportaciones & Una misma fuente cumple funciones distintas según el perfil que la mira; obligar a elegir una sola ubicación traslada al lector una decisión que el propio contenido no tiene resuelta \\
\end{uxtabla}

La única predicción que el documento había dejado escrita antes de aplicar la prueba se refería a ``Credenciales API''. La hipótesis derivada del escenario de Ximena Solís Barrera sostenía que esa tarjeta debía permanecer visible en un segundo nivel y no enterrarse dentro de una jerarquía técnica más profunda. \textbf{La prueba la confirmó}: 7 de los 8 evaluadores resolvieron en Administración, al primer intento, la tarea de crear una credencial de integración, y solo 3 de los 8 titubearon antes de decidir; el único desvío fue hacia Inicio, no hacia una rama técnica más honda. A ello se sumó que 4 de los 8 evaluadores pidieron duplicar esa misma tarjeta, señal de que el contenido se reconoce desde más de un punto de la estructura. La tarjeta se conservó en segundo nivel, dentro de ``4.4 Integraciones'', y se marcó como faceta transversal.

\begin{uxnota}[Dos predicciones falsables, dos desenlaces distintos]
El documento llegó a la prueba con dos apuestas escritas de antemano: que ``Credenciales API'' funcionaría en un segundo nivel y que los tableros funcionarían bajo Gobierno. La primera salió bien, con 7 aciertos de 8, y el mapa la conservó sin tocarla. La segunda salió mal, con 1 acierto de 8, y el mapa cambió. Este entregable reporta los dos desenlaces, el confirmado y el refutado, porque de esa convivencia depende que la validación se pueda auditar.
\end{uxnota}

\subsection{Implementación y Pruebas}
El mapa revisado funciona como el \textbf{contrato de navegación} del producto: sus cuatro módulos y sus subniveles se traducen uno a uno en la estructura de rutas del prototipo, así que cada rama del árbol tendrá una ruta correspondiente y ninguna ruta existirá sin rama que la sustente. Sobre ese contrato se construirá la interfaz de alta fidelidad en la Actividad 4, incluidas la barra lateral con estado activo, la rama 2.4 y el acceso cruzado a la bitácora desde Gobierno del dato. La verificación con personas reales corresponde a la Actividad 5: una prueba de usabilidad con al menos cinco participantes, medida con la escala SUS y meta de 75 puntos o más, acompañada del registro de tasa de éxito y de ruta elegida en las mismas cinco tareas usadas aquí, para comparar el desempeño del árbol revisado contra las 28 coincidencias de 40 obtenidas en esta corrida.

\subsection{Iteración y Mejora Continua}
El mapa de navegación de Karisma Data se mantiene bajo revisión periódica. Establecimos una política de gobernanza ya adoptada por el equipo: la estructura se revisa y actualiza cada trimestre a partir de las analíticas de clics y de las búsquedas fallidas, y toda modificación posterior se documenta en la misma tabla de cambios de cuatro columnas que se usó en el paso 4. La corrida del 9 de agosto de 2026 quedó fijada como línea base de esa política.
